% !TeX root = ../papers/nc-p3.tex
\section{六次式の正則一点が対称性を外した窓全体でも独立な成分になること}
\label{sec:isolated-split-component}
\begin{theorem}\label{thm:isolated-split-component}
定理\ref{thm:split-bell-locus}の六次式から得られる無印の一類を $[w]$ とする。
隣接二切断の階数がともに6以下となる四テンソルの射影スキームを $Y$、
基底群を $G=\prod_{i=0}^3\mathrm{PGL}(V_i)$ とする。
$Y$ は $[w]$ で滑らかで、その唯一の既約成分は $\overline{G[w]}$、次元は59である。
$G[w]$ は $Y$ の開集合で、この成分の半安定部分全体に等しい。
従って半安定商ではこの最小周期1の左右AS正則窓は被約な開かつ閉な一点である。
標準窓の任意の族で $[w]$ に特殊化する幾何学的一般ファイバーはこれと添字保持同型である。
\end{theorem}
\begin{proof}[標数零の代表と良い剰余]
分解の八つの二文字空間を定理\ref{thm:split-inverse-transporter}の順とする。
$P$ のBell四数は定理\ref{thm:split-bell-locus}第三型の $(1,b,c,d)$ で、
$(d^2+1)^3-4d^3=0$ とする。$\iota^2=-1$ とし、第一切断の四対応を恒等、
第二切断の四対応を順に
\[
 \begin{pmatrix}1&1\\-\iota&\iota\end{pmatrix},\quad
 \tfrac12\begin{pmatrix}1&1\\\iota&-\iota\end{pmatrix},\quad
 \begin{pmatrix}1&-1\\\iota&\iota\end{pmatrix},\quad
 \tfrac12\begin{pmatrix}1&-1\\-\iota&-\iota\end{pmatrix}
\]
に取る。両逆転置の余因子をこれらの対応で運び、第二の組と第一の組の射影比例を
六次式を法として全十六座標で検算した。従って六ブロックのSchur式の倍率が決まり、
同定理により標数零の代表となる。$109$ を法とし $d=17,\iota=33$ に取れる。
六次式の根は単純で、両切断行列式、倍率の規格化に用いた座標、四対応の行列式は
全て単元である。この整数環の局所化上で両方の切断を比較できる。
\end{proof}
\begin{proof}[制約なしの接空間]
$16\times16$ の階数6行列 $M$ において七次小行列の一次条件は、
左右の十次元核の基底を行列 $L,R$ として $L\dot M R=0$ である。
二切断それぞれから100行を得て、全256係数を列とする。重みは法線文字の個数に取る。
基点の重みが一定なので核は重みごとに分解する。
前記の良い剰余での列数、制約階数、接次元は順に
\[
\begin{array}{c|rrrrr}
\text{法線文字数}&0&1&2&3&4\\\hline
\text{列数}&16&64&96&64&16\\
\text{階数}&16&48&68&48&16\\
\text{核次元}&0&16&28&16&0 .
\end{array}
\]
全行列の階数も196である。\path{computations/verify_isolated_split_component.py}
は上述の体上の射影比例、剰余での四テンソルの構成、両核からの全行列を再現する。
六次小行列の単元に沿って両核の基底を持ち上げられるので、標数零のアフィン接次元は60以下となる。
一方、定理\ref{thm:split-bell-locus}により射影安定群は一次元トーラス、
射影軌道の次元は $4(4^2-1)-1=59$ である。原点以外の錐でスカラーを戻すと60となる。
従って接次元は局所次元と一致し、標数零で滑らかである。
その点を通る成分はただ一つで、その次元に等しい軌道の閉包がこの成分になる。
軌道はその閉包で開、他成分はそのどの点も含まないので $Y$ でも開である。

この軌道は半安定部分で閉であることも定理\ref{thm:split-bell-locus}で示した。
従って軌道閉包のそれ以外の点は全て不安定である。商でこの軌道と同じ像を持つ
半安定点は、自身の軌道閉包がこの開軌道に会うので初めから同じ軌道に属する。
従って半安定部分の開かつ閉な飽和集合であり、その商は被約一点である。
最後の族に関する主張も軌道の開性と二窓の同型定理による。
\end{proof}
